% ============================================================
% DRAFT NOTE (target)
% Contains: The setting without notation; the paired damage event and its
% three consequences; the per-step risk target; risk vs magnitude; the
% current-state boundary; earliness defined via absolute checkpoints.
% Written now: Complete per docs/goal.md and quality_risk_protocol.md
% sections 1 and 4.
% Pending: Nothing result-dependent. Task scorers are named in the corpus
% section, not here.
% Sources: docs/goal.md; docs/implementation/quality_risk_protocol.md
% ============================================================
\section{Damage and the Risk Target}
\label{sec:target}

Consider a model decoding an answer with its KV cache under a known
compressor at a known removal ratio, and consider the answer the same model
would have produced for the same prompt with its cache intact. If the
compressed answer is worse at the task, compression did harm; if it is as
good or better, it did not, whatever else changed along the way. We want to
know, while the compressed answer is still being written, how likely the
first outcome is. This section defines that quantity.

\paragraph{The damage event.}
Fix a prompt \prompt{} and a compression action $\action = (\press, \ratio)$, a
compressor \press{} at removal ratio \ratio. Let \yfull{} be the reference
run, generated for \prompt{} with the full cache, and \ycmp{} the compressed
run, generated for \prompt{} under \action{} from the first decoded token
onward. Let $\score(\cdot)$ be the task's own score for a completed output. The compressed run is \emph{damaged} when
\begin{equation}
\dmg \;=\; \ind{\, \score(\yfull) - \score(\ycmp) > 0 \,}.
\label{eq:damage}
\end{equation}
where $\ind{\cdot}$ is the indicator function, equal to one when the
condition inside holds and zero otherwise; \dmg{} is therefore a binary
label, one per compressed run. Three consequences follow, and each is
deliberate. A compressed run that is
wrong when the reference run is also wrong is not damaged: the model would
have failed anyway, and charging that failure to compression would credit the
forecaster for predicting the model's baseline error rate. A compressed run
that scores higher than the reference is not damaged; compression-associated
lift exists and is not harm. And a compressed run that produces different
tokens, a different distribution, or different wording, yet scores the same,
is not damaged. Damage is defined in task space, against the paired
counterfactual, and nowhere else.

\paragraph{The risk target.}
The forecast is made while \ycmp{} is being generated. At decoding
step $t$, with $t$ tokens of the compressed run emitted, the target is
\begin{equation}
\risk(\step) \;=\; P\!\left(\dmg = 1 \,\middle|\, \info\right),
\label{eq:risk}
\end{equation}
where \info{} is the information causally available at \step: the
compression action, the position, and statistics of the compressed model's
own next-token distributions up to and including step $t$. The reference
run, its score, and anything downstream of $t$ are not in \info.
We call $\risk$ the hazard: the current estimated risk that the run ends
damaged. The label is one per run and does not change with \step; what changes is the
evidence. Every token of a run is a training row carrying the run's \dmg, and
the forecaster's job is to move $\risk(\step)$ toward the right answer as early
as the evidence allows. When no compression is active, \dmg{} is zero by
definition and $\risk(\step) \equiv 0$ without invoking any model; this anchors
the scale and lets one forecaster run continuously across compressed and
uncompressed operation.

\paragraph{Risk, not magnitude.}
We forecast the probability of damage rather than the size of
$\score(\yfull) - \score(\ycmp)$. Operationally, a probability is
what a downstream decision needs: whether to back off depends on comparing
the risk against the cost of backing off, not on a point estimate of lost
score. Statistically, a binary event on labels that arrive once per run is
a far more learnable target than a signed magnitude on the same labels.

\paragraph{Current-state boundary.}
Equation~\ref{eq:risk} conditions on compression that is already active and
stays active for the rest of the run. It does not answer the prospective
question, what would happen if compression were switched on now for a run
decoding with a full cache, nor the question of changing the ratio mid-run.
Those are counterfactual interventions with their own label requirements and
are out of scope.

\paragraph{Early enough.}
Because the label is per run and the forecast per token, the useful summary
of a forecaster is not a single number but a curve: how well $\risk(\step)$
separates damaged from undamaged runs, and how well calibrated it is, as a
function of $t$. We evaluate at fixed absolute checkpoints
$t \in \{8, 16, 32, 64, 128\}$. A run contributes to a checkpoint only if it
is still generating at that position; runs that have already finished are
excluded there rather than carried forward. Absolute positions are the only
admissible clock. Relative progress, the fraction of the run completed, is
computed from the final length and therefore leaks the outcome: runs that
loop or fail to terminate are long, and a forecaster given their relative
progress would be told, in effect, how they end.
